% Unit 074 terjemahan Bahasa Indonesia dari chapter6.tex baris 557--739.
% Sumber: Methods of Algebra, Volume 2, commit
% 9a5803ff77dd3257484cb177f851a73770a59dd3 (CC BY 4.0).
% stable-unit-id: o014.aljabr2.chapter6.low-degree-cohomology
% Cakupan: Bagian 6.3 lengkap.

% segment-id: o014.li2.chapter6.low-degree-cohomology.h001
\section{Kohomologi berderajat rendah: homomorfisme silang dan ekstensi grup}
\label{sec:grp-low-degree}
\hypertarget{o014.aljabr2.chapter6.low-degree-cohomology}{ }

% segment-id: o014.li2.chapter6.low-degree-cohomology.p001
Bagian ini bertujuan memperdalam pemahaman tentang $\Hm^1$ dan $\Hm^2$.
Misalkan $G$ grup dan $\Bbbk$ gelanggang komutatif. Untuk modul-$G$ $M$,
untuk kompleks standar $C(G,M)$ dalam Proposisi~\ref{prop:groupcoh-cochain},
definisikan submodul-submodul dari $C^n(G,M)$
\[ Z^n(G,M):=\Ker(d^n)\supset\Image(d^{n-1})=:B^n(G,M), \]
dan definisikan pula $B^0(G,M)=0$. Kita mengidentifikasi
\begin{equation*}
	\Hm^n(G,M)\quad\text{dengan}\quad Z^n(G,M)/B^n(G,M).
\end{equation*}

% segment-id: o014.li2.chapter6.low-degree-cohomology.p002
Karena asal-usulnya dalam topologi, unsur $Z^n(G,M)$ juga disebut
$n$-\emph{kosiklus}, sedangkan unsur $B^n(G,M)$ disebut
$n$-\emph{kobatas}.
\index{kosiklus (cocycle)}
\index{kobatas (coboundary)}
\index[sym1]{ZGM@$Z^n(G,M)$}
\index[sym1]{BGM@$B^n(G,M)$}

% segment-id: o014.li2.chapter6.low-degree-cohomology.p003
Ingat bahwa $C^n(G,M)$ terdiri atas semua pemetaan $f:G^n\to M$. Bagian
berderajat rendah dari kompleks standar adalah
\[\begin{tikzcd}[row sep=tiny, column sep=small]
	C^0(G,M) \arrow[r,"{d^0}"] & C^1(G,M) \arrow[r,"{d^1}"] & C^2(G,M) \\
	m \arrow[mapsto,r] & {[g\mapsto gm-m]} & \\
	& f \arrow[mapsto,r] &
	{[(g_1,g_2)\mapsto g_1f(g_2)-f(g_1g_2)+f(g_1)]},
\end{tikzcd}\]
Khususnya, $Z^0(G,M)=M^G$. Berikutnya kita meninjau kasus $n=1$.

% segment-id: o014.li2.chapter6.low-degree-cohomology.d001
\begin{definition}[Homomorfisme silang]\label{def:crossed-homomorphism}
	\index{homomorfisme silang (crossed homomorphism)}
	Misalkan $M$ modul-$G$. Unsur-unsur $Z^1(G,M)$ juga disebut
	\emph{homomorfisme silang} dari $G$ ke $M$; dengan kata lain,
	homomorfisme silang adalah pemetaan $f:G\to M$ yang memenuhi
	\[ f(g_1g_2)=g_1f(g_2)+f(g_1),\qquad g_1,g_2\in G. \]
	Untuk $\Bbbk$ tetap, homomorfisme-homomorfisme ini membentuk
	modul-$\Bbbk$ terhadap operasi titik demi titik.
\end{definition}

% segment-id: o014.li2.chapter6.low-degree-cohomology.p004
Perhatikan bahwa unsur $B^1(G,M)$ adalah homomorfisme silang berbentuk
$f_m:g\mapsto gm-m$, dengan $m\in M$.

% segment-id: o014.li2.chapter6.low-degree-cohomology.p005
Secara umum, homomorfisme silang $f$ selalu memenuhi $f(1_G)=0$ (ambil
$g_1=g_2=1_G$) dan $g^{-1}f(g)=-f(g^{-1})$ (ambil $g_1=g^{-1}$ dan $g_2=g$).

% segment-id: o014.li2.chapter6.low-degree-cohomology.e001
\begin{example}
	Jika aksi $G$ pada $M$ trivial, maka homomorfisme silang tepat merupakan
	homomorfisme dari $G$ ke grup aditif $(M,+)$, dan $B^1(G,M)=\{0\}$. Jadi
	dalam keadaan ini
	\[ \Hom_{\cate{Grp}}(G,M)\simeq\Hm^1(G,M). \]
\end{example}

% segment-id: o014.li2.chapter6.low-degree-cohomology.p006
Cara lain memahami homomorfisme silang adalah melalui hasil kali semilangsung.
Untuk modul-$G$ $M$, gunakan aksi $G$ untuk membangun hasil kali semilangsung
$M\rtimes G$. Unsurnya berbentuk $(m,g)$ dengan $m\in M$, $g\in G$, dan
terdapat homomorfisme proyeksi $\pi:M\rtimes G\to G$ yang memetakan $(m,g)$
ke $g$.

% segment-id: o014.li2.chapter6.low-degree-cohomology.t001
\begin{proposition}\label{prop:crossed-semidirect-product}
	Misalkan $M$ modul-$G$. Terdapat bijeksi kanonik
	\[ Z^1(G,M)\xrightarrow{1:1}
	\left\{\text{homomorfisme grup }\sigma:G\to M\rtimes G
	\;\middle|\;\pi\sigma=\identity_G\right\}. \]
	Pemetaan $\sigma$ di ruas kanan juga disebut penampang dari $\pi$;
	homomorfisme silang $f$ bersesuaian dengan penampang
	$\sigma(g)=(f(g),g)$.
\end{proposition}
\begin{proof}
	Pemetaan $\sigma$ yang memenuhi $\pi\sigma=\identity_G$ harus berbentuk
	$\sigma(g)=(f(g),g)$ untuk suatu pemetaan $f:G\to M$. Berdasarkan definisi
	hasil kali semilangsung, $\sigma$ merupakan homomorfisme grup jika dan hanya
	jika dalam $M$ berlaku
	\[ f(g_1)+g_1f(g_2)=f(g_1g_2),\qquad g_1,g_2\in G. \]
	Ini tepat definisi homomorfisme silang.
\end{proof}

% segment-id: o014.li2.chapter6.low-degree-cohomology.p007
Homomorfisme silang juga berkaitan dengan ideal augmentasi
$\mathfrak{I}\subset\Bbbk[G]$ yang diperkenalkan dalam
Catatan~\ref{rem:augmentation-kG}.

% segment-id: o014.li2.chapter6.low-degree-cohomology.t002
\begin{proposition}\label{prop:crossed-augmentation}
	Untuk setiap modul-$G$ $M$, terdapat isomorfisme kanonik
	$\Hom_G(\mathfrak{I},M)\rightiso Z^1(G,M)$ yang memetakan homomorfisme
	$\varphi$ ke homomorfisme silang $f(g)=\varphi(g-1_G)$.
\end{proposition}
\begin{proof}
	Sebagai modul-$\Bbbk$, $\mathfrak{I}$ mempunyai basis
	$\{g-1_G:g\in G,\;g\neq1_G\}$. Jadi menentukan homomorfisme modul-$\Bbbk$
	$\varphi:\mathfrak{I}\to M$ setara dengan menentukan pemetaan $f:G\to M$
	yang memenuhi $f(1_G)=0$, dengan hubungan
	$f(g)=\varphi(g-1_G)$. Syarat agar $\varphi$ merupakan homomorfisme
	modul-$G$ adalah
	$\varphi(g_1(g_2-1_G))=g_1\varphi(g_2-1_G)$. Akan tetapi,
	\[ \text{ruas kiri}=\varphi(g_1g_2-1_G-g_1+1_G)
	=f(g_1g_2)-f(g_1),\qquad
	\text{ruas kanan}=g_1f(g_2). \]
	Ini tepat syarat homomorfisme silang.
\end{proof}

% segment-id: o014.li2.chapter6.low-degree-cohomology.p008
Untuk langkah berikutnya, yaitu kasus $n=2$, kita harus mengingat kembali
konsep ekstensi grup. Ambil $\Bbbk=\Z$ dan perkenalkan beberapa istilah.

% segment-id: o014.li2.chapter6.low-degree-cohomology.d002
\begin{definition}\label{def:group-ext}
	\index{ekstensi grup (extension of groups)}
	Misalkan $A$ grup. Sebuah \emph{ekstensi} $G$ oleh $A$ adalah barisan eksak
	pendek grup \cite[Definisi 4.3.7]{Li1}
	\[ 0\to A\to E\xrightarrow{\pi}G\to1. \]

	Selanjutnya ekstensi grup ini cukup dilambangkan dengan $E$. Homomorfisme grup
	$s:G\to E$ yang memenuhi $\pi s=\identity_G$ disebut suatu
	\emph{pembelahan} ekstensi tersebut. Dalam keadaan ini, $s$ menyertakan $G$
	sebagai subgrup dari $E$, dan melalui penyertaan ini $E$ diidentifikasi
	dengan hasil kali semilangsung $A\rtimes G$. Ekstensi yang mempunyai
	pembelahan disebut \emph{terbelah}. Untuk rinciannya, lihat
	\cite[\S 4.3]{Li1}.
\end{definition}

% segment-id: o014.li2.chapter6.low-degree-cohomology.p009
Secara garis besar, teori ekstensi grup mengkaji cara membangun grup yang lebih
besar dari subgrup normal dan grup hasil bagi. Untuk menghubungkannya dengan
kohomologi modul-$G$, selanjutnya grup $A$ di atas ditulis sebagai $M$ dan
disyaratkan komutatif, dengan operasi grup ditulis secara aditif. Ekstensi
grup semacam ini menentukan homomorfisme grup
\[ \alpha:G\to\Aut_{\cate{Grp}}(M),\qquad
	\alpha(g)(x)=exe^{-1},\qquad e\in\pi^{-1}(g)\ \text{dipilih sembarang}. \]
Dengan demikian $M$ menjadi modul-$G$ melalui $\alpha$; ini merupakan
invarian paling dasar dari ekstensi grup.

% segment-id: o014.li2.chapter6.low-degree-cohomology.p010
Untuk modul-$G$ $M$ yang diberikan, kita ingin mengklasifikasikan semua
ekstensi $G$ oleh $M$ yang $\alpha$-nya tepat merupakan struktur modul-$G$
pada $M$, hingga ekuivalensi (atau isomorfisme) ekstensi. Ekuivalensi dari
ekstensi $E$ ke $E'$ didefinisikan secara alami sebagai diagram komutatif grup
\begin{equation}\label{eqn:group-ext-diagram}\begin{tikzcd}
	0 \arrow[r] & M \arrow[r] \arrow[equal,d] &
	E \arrow[r] \arrow[d,"\varphi"] & G \arrow[r] \arrow[equal,d] & 1 \\
	0 \arrow[r] & M \arrow[r] & E' \arrow[r] & G \arrow[r] & 1.
\end{tikzcd}\end{equation}
Pembaca diminta memverifikasi bahwa $\varphi$ dalam diagram tersebut secara
otomatis merupakan isomorfisme grup, dan bahwa suatu ekstensi terbelah jika
dan hanya jika ekuivalen dengan hasil kali semilangsung $M\rtimes G$.

% segment-id: o014.li2.chapter6.low-degree-cohomology.d003
\begin{definition}
	\index[sym1]{ExtGM@$\cate{Ext}(G,M)$}
	Definisikan kategori $\cate{Ext}(G,M)$ dari semua ekstensi $G$ oleh modul-$G$
	$M$ yang diberikan; morfismenya adalah diagram komutatif
	\eqref{eqn:group-ext-diagram}.
\end{definition}

% segment-id: o014.li2.chapter6.low-degree-cohomology.p011
Kategori ini merupakan grupoid: semua morfismenya dapat dibalik.
Mengklasifikasikan ekstensi berarti mendeskripsikan $\cate{Ext}(G,M)$ hingga
ekuivalensi kategori. Persoalan ini terbagi menjadi dua:
\begin{itemize}
	\item mendeskripsikan $\Obj(\cate{Ext}(G,M))/\simeq$, yakni kelas
	ekuivalensi ekstensi, dan mengenali ekstensi terbelah di antaranya;
	\item mendeskripsikan morfisme dalam $\cate{Ext}(G,M)$, yakni isomorfisme
	antara ekstensi.
\end{itemize}
Kita mulai dari persoalan kedua yang lebih sederhana. Mula-mula, suatu
ekstensi $E$ mempunyai kelas automorfisme sederhana berbentuk
$\Ad_m:e\mapsto mem^{-1}$, dengan $m\in M$ dan operasi grup dalam $E$ ditulis
secara multiplikatif; automorfisme ini disebut automorfisme dalam. Karena
alasan abstrak, automorfisme dalam membentuk subgrup normal dari $\Aut(E)$,
dan grup hasil bagi terkait disebut grup automorfisme luar dari $E$.

% segment-id: o014.li2.chapter6.low-degree-cohomology.p012
Untuk mengklasifikasikan ekstensi grup, kita perlu memakai kompleks standar
ternormalisasi $\overline{C}(G,M)$ yang dijelaskan dalam
Catatan~\ref{rem:nor-cochain} untuk menghitung kohomologi. Definisikan
\begin{gather*}
	\overline{C}^n(G,M)\supset\overline{Z}^n(G,M)
	\supset\overline{B}^n(G,M),\\
	\Hm^n(G,M)\simeq\overline{Z}^n(G,M)/\overline{B}^n(G,M).
\end{gather*}
Kita mempunyai $\overline{C}^0(G,M)=C^0(G,M)=M$. Pembahasan sebelumnya
tentang homomorfisme silang juga telah menunjukkan bahwa
$\overline{Z}^1(G,M)=Z^1(G,M)$.

% segment-id: o014.li2.chapter6.low-degree-cohomology.l001
\begin{lemma}\label{prop:factor-set-prep}
	Misalkan $0\to M\to E\xrightarrow{\pi}G\to1$ merupakan objek
	$\cate{Ext}(G,M)$. Grup automorfismenya secara kanonik isomorfik dengan
	$\overline{Z}^1(G,M)$, sedangkan grup automorfisme luarnya secara kanonik
	isomorfik dengan $\Hm^1(G,M)$. Lebih eksplisit, unsur
	$z\in\overline{Z}^1(G,M)$ bersesuaian dengan automorfisme yang memetakan
	$e\in E$ ke $z(\pi(e))e$.
\end{lemma}
\begin{proof}
	Untuk sementara tuliskan operasi grup pada $M$ secara multiplikatif. Ambil
	$E'=E$ dalam diagram komutatif \eqref{eqn:group-ext-diagram}. Untuk setiap
	$e\in E$ terdapat $z(e)\in M$ tunggal dengan $\varphi(e)=z(e)e$. Jika
	$m\in M$, maka
	$\varphi(me)=m\varphi(e)=mz(e)e=z(e)me$, sehingga $z(e)$ hanya bergantung
	pada citra $e$ dalam $G$; selanjutnya kita menulisnya sebagai pemetaan
	$z:G\to M$. Sebaliknya, setiap pemetaan $z:G\to M$ menentukan pemetaan
	$\varphi:E\to E$ melalui $\varphi(e)=z(\pi(e))e$, dan diagram
	\eqref{eqn:group-ext-diagram} komutatif pada tingkat himpunan.

	Perhatikan pemetaan $z:G\to M$ dan $\varphi$ terkait. Misalkan
	$g_i=\pi(e_i)$ untuk $i=1,2$. Maka
	\begin{align*}
		\varphi(e_1e_2)&=z(g_1g_2)e_1e_2,\\
		\varphi(e_1)\varphi(e_2)&=z(g_1)e_1z(g_2)e_2\\
		&=z(g_1)\underbracket{g_1\cdot z(g_2)}_{
		\text{struktur modul-$G$ pada $M$}}e_1e_2.
	\end{align*}
	Jadi $\varphi$ merupakan homomorfisme grup jika dan hanya jika $z$ adalah
	homomorfisme silang. Jelas bahwa komposisi homomorfisme grup bersesuaian
	dengan penjumlahan homomorfisme silang. Sekarang ambil $m\in M$.
	Automorfisme dalam terkait ialah
	$e\mapsto mem^{-1}=mem^{-1}e^{-1}e$. Jelas bahwa
	$mem^{-1}e^{-1}\in M$, dan dalam notasi aditif unsur ini adalah
	$m-\pi(e)m$. Selesai.
\end{proof}

% segment-id: o014.li2.chapter6.low-degree-cohomology.d004
\begin{definition}
	Grupoid $\tau^{\leq2}\overline{\cate{C}}(G,M)$ didefinisikan sebagai
	berikut: himpunan objeknya adalah $\overline{Z}^2(G,M)$, sedangkan himpunan
	morfismenya adalah
	\[ \Hom(f,f'):=\{w\in\overline{C}^1(G,M):f'=d^1w+f\}. \]
	Komposisi morfisme didefinisikan dengan penjumlahan dalam
	$\overline{C}^1(G,M)$. Dengan demikian, kelas isomorfisme objek berkorespondensi
	secara bijektif dengan unsur $\Hm^2(G,M)$.
\end{definition}

% segment-id: o014.li2.chapter6.low-degree-cohomology.p013
Makna notasinya jelas: $\tau^{\leq2}\overline{\cate{C}}(G,M)$ adalah
kategori yang ditentukan oleh kompleks terpotong
$\tau^{\leq2}\overline{C}(G,M)$.

% segment-id: o014.li2.chapter6.low-degree-cohomology.t003
\begin{theorem}\label{prop:factor-set}
	Misalkan $M$ modul-$G$. Kategori $\cate{Ext}(G,M)$ ekuivalen dengan
	$\tau^{\leq2}\overline{\cate{C}}(G,M)$. Khususnya, kelas ekuivalensi
	ekstensi $G$ oleh $M$ berkorespondensi satu-satu dengan unsur
	$\Hm^2(G,M)$, dan kelas ekuivalensi ekstensi terbelah bersesuaian dengan
	$0\in\Hm^2(G,M)$.
\end{theorem}
\begin{proof}
	Mula-mula kita bangun funktor
	$\tau^{\leq2}\overline{\cate{C}}(G,M)\to\cate{Ext}(G,M)$. Ambil
	pemetaan $f:G^2\to M$ dan definisikan operasi biner pada himpunan $M\times G$
	dengan
	\[ (m_1,g_1)\cdot(m_2,g_2):=
	\left(m_1+g_1\cdot m_2+f(g_1,g_2),g_1g_2\right). \]
	Sifat asosiatif ekuivalen dengan identitas
	\begin{equation}\label{eqn:grp-2-cocycle}
		f(g_1,g_2)+f(g_1g_2,g_3)
		=g_1\cdot f(g_2,g_3)+f(g_1,g_2g_3),
	\end{equation}
	yakni $f\in Z^2(G,M)$. Lebih lanjut, $(0,1_G)$ merupakan unsur satuan jika
	dan hanya jika
	\[ f(g,1_G)=0=f(1_G,g), \]
	yakni $f\in\overline{Z}^2(G,M)$. Dengan mengambil
	$(g_1,g_2,g_3)=(g,g^{-1},g)$ dalam \eqref{eqn:grp-2-cocycle} dan memakai
	persamaan terakhir, diperoleh
	\begin{equation}\label{eqn:grp-2-cocycle-inv}
		f(g,g^{-1})=g\cdot f(g^{-1},g).
	\end{equation}

	Jadi $M\times G$ memperoleh struktur monoid dari
	$f\in\overline{Z}^2(G,M)$; tuliskan monoid tersebut sebagai $E$. Ia secara
	otomatis merupakan grup, sebab dari \eqref{eqn:grp-2-cocycle-inv} diperoleh
	\[ (0,g)^{-1}=\left(-f(g^{-1},g),g^{-1}\right),\qquad
	(m,1_G)^{-1}=(-m,1_G). \]
	Mudah dibuktikan bahwa pemetaan $m\mapsto(m,1_G)$ dan $(m,g)\mapsto g$
	memberikan ekstensi grup $0\to M\to E\to G\to1$ yang menginduksi tepat
	aksi $G$ semula pada $M$. Periksalah bahwa $f=0$ bersesuaian dengan
	hasil kali semilangsung $E=M\rtimes G$.

	Jika selain itu $w\in\overline{C}^1(G,M)$ dan $f':=d^1w+f$, tuliskan
	ekstensi terkait sebagai $E'$. Pemetaan
	\[ (m,g)\mapsto(m-w(g),g) \]
	memberikan isomorfisme dari $E$ ke $E'$. Memang, syarat bahwa pemetaan ini
	mempertahankan perkalian setara dengan identitas
	\[ -w(g_1)-g_1\cdot w(g_2)+(d^1w)(g_1,g_2)=-w(g_1g_2), \]
	yang tepat merupakan definisi $d^1w$.

	Dengan demikian diperoleh funktor
	$\tau^{\leq2}\overline{\cate{C}}(G,M)\to\cate{Ext}(G,M)$, dan kedua
	kategori merupakan grupoid. Dengan membandingkan grup automorfisme objek
	melalui Lema~\ref{prop:factor-set-prep}, funktor ini penuh dan setia. Kini
	kita tunjukkan bahwa funktor tersebut surjektif esensial.

	Misalkan $E$ objek dalam $\cate{Ext}(G,M)$. Pilih penampang sembarang
	$s:G\to E$ dari pemetaan $\pi:E\to G$ dalam data ekstensi, yaitu pemetaan
	yang memenuhi $\pi s=\identity_G$. Ini setara dengan memilih wakil dalam
	setiap koset $M$ di $E$. Penampang $s$ disebut ternormalisasi jika
	$s(1_G)=1_E$. Untuk semua $g_1,g_2\in G$, unsur $s(g_1g_2)$ dan
	$s(g_1)s(g_2)$ mempunyai citra yang sama di bawah $\pi$. Jadi terdapat
	$f(g_1,g_2)\in M$ tunggal dengan
	\[ s(g_1)s(g_2)=f(g_1,g_2)s(g_1g_2). \]

	Mudah dibuktikan bahwa sifat asosiatif
	$(s(g_1)s(g_2))s(g_3)=s(g_1)(s(g_2)s(g_3))$ diterjemahkan menjadi
	\eqref{eqn:grp-2-cocycle}, yakni $f\in Z^2(G,M)$. Jika $s$ ternormalisasi,
	maka $f(1_G,g)=0=f(g,1_G)$ untuk semua $g$, yang setara dengan
	$f\in\overline{Z}^2(G,M)$.

	Untuk penampang ternormalisasi $s$ yang dipilih, beri $M\times G$ struktur
	grup yang bersesuaian dengan $f$. Perhatikan bijeksi
	\[ M\times G\xrightarrow{1:1}E,\qquad(m,g)\mapsto ms(g). \]
	Dari pembahasan di atas, bijeksi ini mempertahankan perkalian grup dan
	memberikan isomorfisme dalam $\cate{Ext}(G,M)$. Ini menyelesaikan bukti
	ekuivalensi.
\end{proof}

% segment-id: o014.li2.chapter6.low-degree-cohomology.p014
Dalam bagian terakhir bukti kesurjektifan esensial, mengubah penampang $s$
setara dengan memilih pemetaan sembarang $h:G\to M$ dan mengambil
$s'(g):=h(g)s(g)$. Pemetaan $f':G^2\to M$ terkait berubah menjadi
\begin{align*}
	f'(g_1,g_2)&=\left(h(g_1)+g_1\cdot h(g_2)-h(g_1g_2)\right)+f(g_1,g_2)\\
	&=(d^1h)(g_1,g_2)+f(g_1,g_2).
\end{align*}
Jika disyaratkan pula $h(1_G)=0$, maka
$d^1h\in\overline{B}^2(G,M)$. Secara historis, pengamatan ini
merupakan titik awal hubungan antara ekstensi dan $\Hm^2(G,M)$, sekaligus
memberikan interpretasi melalui ekstensi grup bagi definisi
$\overline{Z}^2(G,M)$ dan $\overline{B}^2(G,M)$, meskipun hal itu tidak
ditonjolkan dalam pembuktiannya.

% segment-id: o014.li2.chapter6.low-degree-cohomology.r001
\begin{remark}\label{rem:Ext-incarnation}
	Teorema~\ref{prop:factor-set} mengatakan bahwa kompleks
	$\tau^{\leq2}\overline{C}(G,M)$ merupakan realisasi linear dari kategori
	ekstensi grup $\cate{Ext}(G,M)$: suku berderajat $2$ mendeskripsikan objek,
	suku berderajat $1$ mendeskripsikan morfisme, dan yang masih perlu
	ditafsirkan adalah suku berderajat $0$. Untuk
	$m\in M=\overline{C}^0(G,M)$, terdapat automorfisme dalam ekstensi grup
	$\Ad_m:e\mapsto mem^{-1}$, yang trivial jika dan hanya jika
	$m\in M^G=\Hm^0(G,M)$. Dari sudut pandang automorfisme dalam, kita dapat
	memandang unsur $M$ sebagai 2-morfisme dalam $\cate{Ext}(G,M)$: untuk
	$m\in M$ dan
	$\varphi_1,\varphi_2\in\Hom_{\cate{Ext}(G,M)}(E,E')$,
	\[ \text{tafsirkan diagram sel-2}\;
	\begin{tikzcd}[column sep=large]
		E \arrow[bend left=60,r,"\varphi_1",""' {name=U}]
		\arrow[bend right=60,r,"\varphi_2"',"" {name=D}] &
		E' \arrow[Rightarrow,from=U,to=D,"m"]
	\end{tikzcd}
	\;\text{sebagai}\;\varphi_2=\Ad_m\varphi_1. \]

	Perhatikan bahwa $\Ad_m\varphi_1=\varphi_1\Ad_m$ dan
	$\Ad_{m_1}\Ad_{m_2}=\Ad_{m_1+m_2}$. Jadi $\Hm^0(G,M)$ menjadi
	``grup 2-automorfisme'' dari setiap morfisme $\varphi$. Dalam pengertian
	ini, kompleks $\tau^{\leq2}\overline{C}(G,M)$ dapat disebut realisasi
	linear dari 2-kategori $\cate{Ext}(G,M)$. Untuk definisi 2-kategori, lihat
	\cite[\S 3.5]{Li1}.
\end{remark}

% segment-id: o014.li2.chapter6.low-degree-cohomology.p015
Aspek-aspek lain ekstensi grup akan dibahas dalam
\S\ref{sec:group-ext-revisited}. Sebelumnya, kita perlu memperkenalkan lebih
banyak operasi pada homologi dan kohomologi grup.
